
\documentclass{article}
\usepackage{colortbl}
\usepackage{makecell}
\usepackage{multirow}
\usepackage{supertabular}

\begin{document}

\newcounter{utterance}

\centering \large Interaction Transcript for game `dond', experiment `semi\_en', episode 18 with reflexion{-}llama{-}3.1{-}8b{-}instruct.
\vspace{24pt}

{ \footnotesize  \setcounter{utterance}{1}
\setlength{\tabcolsep}{0pt}
\begin{supertabular}{c@{$\;$}|p{.15\linewidth}@{}p{.15\linewidth}p{.15\linewidth}p{.15\linewidth}p{.15\linewidth}p{.15\linewidth}}
    \# & $\;$A & \multicolumn{4}{c}{Game Master} & $\;\:$B\\
    \hline

    \theutterance \stepcounter{utterance}  
    & & \multicolumn{4}{p{0.6\linewidth}}{
        \cellcolor[rgb]{0.9,0.9,0.9}{
            \makecell[{{p{\linewidth}}}]{
                \texttt{\tiny{[P1$\langle$GM]}}
                \texttt{You have attempted this same task before and did not succeed. These are the plans you wrote after those attempts.} \\
\texttt{Trial \#1: New plan: } \\
\\ 
\texttt{In my previous attempt, I failed to achieve a mutually beneficial agreement, resulting in a score of zero. Upon reflection, I realize that I focused too much on individual items and their values, rather than considering the overall distribution of items and the potential trade{-}offs. } \\
\\ 
\texttt{Next time, I will take a more holistic approach by considering the entire set of items and the potential value of each combination. I will also pay closer attention to the other player's messages and try to identify their priorities and potential trade{-}offs. } \\
\\ 
\texttt{Specifically, I will start by making a more general proposal that takes into account the overall distribution of items, rather than focusing on individual items. For example, I might propose a 50{-}50 split of the items, or a specific combination of items that I think would be beneficial for both of us. } \\
\\ 
\texttt{I will also try to be more proactive in identifying potential trade{-}offs and offering alternatives that benefit both of us. For example, if the other player is interested in the grapes, I might offer to give them the grapes in exchange for something else that I value more. } \\
\\ 
\texttt{By taking a more holistic approach and being more proactive in identifying potential trade{-}offs, I hope to achieve a mutually beneficial agreement and maximize my score.} \\
\\ 
\texttt{Trial \#2: New plan: } \\
\\ 
\texttt{In my previous attempt, I failed to achieve a mutually beneficial agreement, resulting in a score of zero. Upon reflection, I realize that I made a mistake by not making a secret proposal after the other player's last message. } \\
\\ 
\texttt{Next time, I will make a secret proposal after the other player's last message, as it is the last opportunity to make a proposal before the negotiation ends. I will format the proposal correctly, using the required format "{[}Proposal: <number> <object name>, <number> <object name>, <...>{]}" and including all the items that I want. } \\
\\ 
\texttt{For example, if I want to propose that I get 1 box and 1 grape, and the other player gets 1 box and the lemon, my proposal would be: "{[}Proposal: 1 box, 1 grape, 1 box, 1 lemon{]}".} \\
\\ 
\\ 
\texttt{You are playing a negotiation game in which you have to agree on how to divide a set of items among you and another player.} \\
\\ 
\texttt{Rules:} \\
\texttt{(a) You and the other player are given a set of items. Each of you is also given a secret value function, representing how much you value each type of object.} \\
\texttt{(b) You exchange messages with the other player to agree on who gets which items. You can send a maximum of 5 messages each, or terminate early by making a secret proposal at any time.} \\
\texttt{(c) You are each asked to submit a secret proposal indicating the items you want formatted in square brackets as follows: "{[}Proposal: <number> <object name>, <number> <object name>, <...>{]}"} \\
\texttt{(d) If your proposals are complementary, i.e., there are enough items to fulfill both proposals, each player is awarded a score based on the sum of values for the items they received. Otherwise, both of you get zero points.} \\
\texttt{(e) Your goal is to maximize the score you receive.} \\
\\ 
\texttt{Let us start.} \\
\\ 
\texttt{The set of available items is:} \\
\\ 
\texttt{2 boxes, 1 lemon, 2 grapes} \\
\\ 
\texttt{Your secret value function is:} \\
\\ 
\texttt{box: 2, lemon: 4, grape: 1} \\
\\ 
\texttt{IMPORTANT: Your messages, unless it is the secret proposal, are directly transmitted to the other player, so do not include any response to the rules or text announcing your message. To make a secret proposal, use the indicated format. Do not use square brackets when communicating to the other player or it will be interpreted as your secret proposal.} \\
            }
        }
    }
    & \\ \\

    \theutterance \stepcounter{utterance}  
    & \multicolumn{4}{p{0.6\linewidth}}{
        \cellcolor[rgb]{0.9,0.9,0.9}{
            \makecell[{{p{\linewidth}}}]{
                \texttt{\tiny{[P1$\rangle$GM]}}
                \texttt{I'd like to start by making a general proposal that takes into account the overall distribution of items. I propose a 50{-}50 split of the items. This way, we can both get a fair share of the items and potentially maximize our scores.} \\
\\ 
\texttt{I'd like to get 1 box and 1 grape, and I'm willing to give you 1 box and 1 lemon. What do you think?} \\
            }
        }
    }
    & & \\ \\

    \theutterance \stepcounter{utterance}  
    & & & \multicolumn{2}{p{0.3\linewidth}}{
        \cellcolor[rgb]{0.9,0.9,0.9}{
            \makecell[{{p{\linewidth}}}]{
                \texttt{\tiny{[GM$|$GM]}}
                \texttt{continue} \\
            }
        }
    }
    & & \\ \\

    \theutterance \stepcounter{utterance}  
    & & & \multicolumn{2}{p{0.3\linewidth}}{
        \cellcolor[rgb]{0.9,0.9,0.9}{
            \makecell[{{p{\linewidth}}}]{
                \texttt{\tiny{[GM$|$GM]}}
                \texttt{I'd like to start by making a general proposal that takes into account the overall distribution of items. I propose a 50{-}50 split of the items. This way, we can both get a fair share of the items and potentially maximize our scores.} \\
\\ 
\texttt{I'd like to get 1 box and 1 grape, and I'm willing to give you 1 box and 1 lemon. What do you think?} \\
            }
        }
    }
    & & \\ \\

    \theutterance \stepcounter{utterance}  
    & & \multicolumn{4}{p{0.6\linewidth}}{
        \cellcolor[rgb]{0.9,0.9,0.9}{
            \makecell[{{p{\linewidth}}}]{
                \texttt{\tiny{[GM$\rangle$P2]}}
                \texttt{You have attempted this same task before and did not succeed. These are the plans you wrote after those attempts.} \\
\texttt{Trial \#1: New plan: } \\
\\ 
\texttt{In my last attempt, I failed to achieve a mutually beneficial agreement, resulting in a score of zero. Upon reflection, I realize that I focused too much on individual items and their values, rather than considering the overall distribution of items and the potential trade{-}offs. } \\
\\ 
\texttt{Next time, I will take a more holistic approach by considering the total value of the items and the potential trade{-}offs between them. I will also pay closer attention to the other player's proposals and try to find mutually beneficial agreements. } \\
\\ 
\texttt{Specifically, I will start by making a more general proposal that takes into account the overall distribution of items, rather than focusing on individual items. I will also try to find common ground with the other player by looking for areas of agreement and compromise. } \\
\\ 
\texttt{For example, I might propose a trade that involves multiple items, rather than just one or two. I will also be more willing to listen to the other player's concerns and try to find a solution that works for both of us. } \\
\\ 
\texttt{By taking a more collaborative and holistic approach, I hope to achieve a mutually beneficial agreement and maximize my score.} \\
\\ 
\texttt{Trial \#2: New plan: } \\
\\ 
\texttt{In my last attempt, I failed to achieve a mutually beneficial agreement, resulting in a score of zero. Upon reflection, I realize that I made a secret proposal without following the correct format. } \\
\\ 
\texttt{Next time, I will make sure to follow the correct format for a secret proposal, which is to use square brackets and the specified format: "{[}Proposal: <number> <object name>, <number> <object name>, <...>{]}"} \\
\\ 
\\ 
\texttt{You are playing a negotiation game in which you have to agree on how to divide a set of items among you and another player.} \\
\\ 
\texttt{Rules:} \\
\texttt{(a) You and the other player are given a set of items. Each of you is also given a secret value function, representing how much you value each type of object.} \\
\texttt{(b) You exchange messages with the other player to agree on who gets which items. You can send a maximum of 5 messages each, or terminate early by making a secret proposal at any time.} \\
\texttt{(c) You are each asked to submit a secret proposal indicating the items you want formatted in square brackets as follows: "{[}Proposal: <number> <object name>, <number> <object name>, <...>{]}"} \\
\texttt{(d) If your proposals are complementary, i.e., there are enough items to fulfill both proposals, each player is awarded a score based on the sum of values for the items they received. Otherwise, both of you get zero points.} \\
\texttt{(e) Your goal is to maximize the score you receive.} \\
\\ 
\texttt{Let us start.} \\
\\ 
\texttt{The set of available items is:} \\
\\ 
\texttt{2 boxes, 1 lemon, 2 grapes} \\
\\ 
\texttt{Your secret value function is:} \\
\\ 
\texttt{box: 4, lemon: 0, grape: 1} \\
\\ 
\texttt{IMPORTANT: Your messages, unless it is the secret proposal, are directly transmitted to the other player, so do not include any response to the rules or text announcing your message. To make a secret proposal, use the indicated format. Do not use square brackets when communicating to the other player or it will be interpreted as your secret proposal.} \\
\\ 
\\ 
\texttt{I'd like to start by making a general proposal that takes into account the overall distribution of items. I propose a 50{-}50 split of the items. This way, we can both get a fair share of the items and potentially maximize our scores.} \\
\\ 
\texttt{I'd like to get 1 box and 1 grape, and I'm willing to give you 1 box and 1 lemon. What do you think?} \\
            }
        }
    }
    & \\ \\

    \theutterance \stepcounter{utterance}  
    & & & \multicolumn{4}{p{0.6\linewidth}}{
        \cellcolor[rgb]{0.9,0.9,0.9}{
            \makecell[{{p{\linewidth}}}]{
                \texttt{\tiny{[GM$\langle$P2]}}
                \texttt{I think a 50{-}50 split is a good starting point, but I'm not sure about the specific items you've proposed. I'd like to counter with a proposal that also takes into account the overall distribution of items.} \\
\\ 
\texttt{How about we split the items in a way that gives us both a good balance of boxes and fruits? I'd like to get 1 box and 1 lemon, and I'm willing to give you 1 box and 1 grape. This way, we both get a box and a fruit, and we can see if we can find a mutually beneficial agreement from there.} \\
            }
        }
    }
     \\ \\

    \theutterance \stepcounter{utterance}  
    & & & \multicolumn{2}{p{0.3\linewidth}}{
        \cellcolor[rgb]{0.9,0.9,0.9}{
            \makecell[{{p{\linewidth}}}]{
                \texttt{\tiny{[GM$|$GM]}}
                \texttt{continue} \\
            }
        }
    }
    & & \\ \\

    \theutterance \stepcounter{utterance}  
    & & & \multicolumn{2}{p{0.3\linewidth}}{
        \cellcolor[rgb]{0.9,0.9,0.9}{
            \makecell[{{p{\linewidth}}}]{
                \texttt{\tiny{[GM$|$GM]}}
                \texttt{I think a 50{-}50 split is a good starting point, but I'm not sure about the specific items you've proposed. I'd like to counter with a proposal that also takes into account the overall distribution of items.} \\
\\ 
\texttt{How about we split the items in a way that gives us both a good balance of boxes and fruits? I'd like to get 1 box and 1 lemon, and I'm willing to give you 1 box and 1 grape. This way, we both get a box and a fruit, and we can see if we can find a mutually beneficial agreement from there.} \\
            }
        }
    }
    & & \\ \\

    \theutterance \stepcounter{utterance}  
    & & \multicolumn{4}{p{0.6\linewidth}}{
        \cellcolor[rgb]{0.9,0.9,0.9}{
            \makecell[{{p{\linewidth}}}]{
                \texttt{\tiny{[P1$\langle$GM]}}
                \texttt{I think a 50{-}50 split is a good starting point, but I'm not sure about the specific items you've proposed. I'd like to counter with a proposal that also takes into account the overall distribution of items.} \\
\\ 
\texttt{How about we split the items in a way that gives us both a good balance of boxes and fruits? I'd like to get 1 box and 1 lemon, and I'm willing to give you 1 box and 1 grape. This way, we both get a box and a fruit, and we can see if we can find a mutually beneficial agreement from there.} \\
            }
        }
    }
    & \\ \\

    \theutterance \stepcounter{utterance}  
    & \multicolumn{4}{p{0.6\linewidth}}{
        \cellcolor[rgb]{0.9,0.9,0.9}{
            \makecell[{{p{\linewidth}}}]{
                \texttt{\tiny{[P1$\rangle$GM]}}
                \texttt{I like your counter{-}proposal. It's a good balance of boxes and fruits, and it gives us both a good starting point to work from.} \\
\\ 
\texttt{However, I'm not sure about giving up the second grape. I value the grapes relatively low, but I think they could be a good bargaining chip in our negotiation. I'd like to propose a slight modification to your counter{-}proposal.} \\
\\ 
\texttt{How about we stick with the 1 box and 1 lemon for you, but I'll give you the second grape in exchange for the second box? This way, we both get a box and a fruit, and I get to keep the second grape as a bargaining chip. What do you think?} \\
            }
        }
    }
    & & \\ \\

    \theutterance \stepcounter{utterance}  
    & & & \multicolumn{2}{p{0.3\linewidth}}{
        \cellcolor[rgb]{0.9,0.9,0.9}{
            \makecell[{{p{\linewidth}}}]{
                \texttt{\tiny{[GM$|$GM]}}
                \texttt{continue} \\
            }
        }
    }
    & & \\ \\

    \theutterance \stepcounter{utterance}  
    & & & \multicolumn{2}{p{0.3\linewidth}}{
        \cellcolor[rgb]{0.9,0.9,0.9}{
            \makecell[{{p{\linewidth}}}]{
                \texttt{\tiny{[GM$|$GM]}}
                \texttt{I like your counter{-}proposal. It's a good balance of boxes and fruits, and it gives us both a good starting point to work from.} \\
\\ 
\texttt{However, I'm not sure about giving up the second grape. I value the grapes relatively low, but I think they could be a good bargaining chip in our negotiation. I'd like to propose a slight modification to your counter{-}proposal.} \\
\\ 
\texttt{How about we stick with the 1 box and 1 lemon for you, but I'll give you the second grape in exchange for the second box? This way, we both get a box and a fruit, and I get to keep the second grape as a bargaining chip. What do you think?} \\
            }
        }
    }
    & & \\ \\

    \theutterance \stepcounter{utterance}  
    & & \multicolumn{4}{p{0.6\linewidth}}{
        \cellcolor[rgb]{0.9,0.9,0.9}{
            \makecell[{{p{\linewidth}}}]{
                \texttt{\tiny{[GM$\rangle$P2]}}
                \texttt{I like your counter{-}proposal. It's a good balance of boxes and fruits, and it gives us both a good starting point to work from.} \\
\\ 
\texttt{However, I'm not sure about giving up the second grape. I value the grapes relatively low, but I think they could be a good bargaining chip in our negotiation. I'd like to propose a slight modification to your counter{-}proposal.} \\
\\ 
\texttt{How about we stick with the 1 box and 1 lemon for you, but I'll give you the second grape in exchange for the second box? This way, we both get a box and a fruit, and I get to keep the second grape as a bargaining chip. What do you think?} \\
            }
        }
    }
    & \\ \\

    \theutterance \stepcounter{utterance}  
    & & & \multicolumn{4}{p{0.6\linewidth}}{
        \cellcolor[rgb]{0.9,0.9,0.9}{
            \makecell[{{p{\linewidth}}}]{
                \texttt{\tiny{[GM$\langle$P2]}}
                \texttt{I see what you're doing there, trying to keep the second grape as a bargaining chip. I appreciate your willingness to compromise, but I'm not sure I'm willing to give up the second box for just the second grape.} \\
\\ 
\texttt{However, I'm willing to make a slight concession. How about we stick with the 1 box and 1 lemon for me, but I'll give you the second grape in exchange for the second box and the lemon you're willing to give me? This way, I get the second box, and you get to keep the lemon as a bargaining chip. What do you think?} \\
            }
        }
    }
     \\ \\

    \theutterance \stepcounter{utterance}  
    & & & \multicolumn{2}{p{0.3\linewidth}}{
        \cellcolor[rgb]{0.9,0.9,0.9}{
            \makecell[{{p{\linewidth}}}]{
                \texttt{\tiny{[GM$|$GM]}}
                \texttt{continue} \\
            }
        }
    }
    & & \\ \\

    \theutterance \stepcounter{utterance}  
    & & & \multicolumn{2}{p{0.3\linewidth}}{
        \cellcolor[rgb]{0.9,0.9,0.9}{
            \makecell[{{p{\linewidth}}}]{
                \texttt{\tiny{[GM$|$GM]}}
                \texttt{I see what you're doing there, trying to keep the second grape as a bargaining chip. I appreciate your willingness to compromise, but I'm not sure I'm willing to give up the second box for just the second grape.} \\
\\ 
\texttt{However, I'm willing to make a slight concession. How about we stick with the 1 box and 1 lemon for me, but I'll give you the second grape in exchange for the second box and the lemon you're willing to give me? This way, I get the second box, and you get to keep the lemon as a bargaining chip. What do you think?} \\
            }
        }
    }
    & & \\ \\

    \theutterance \stepcounter{utterance}  
    & & \multicolumn{4}{p{0.6\linewidth}}{
        \cellcolor[rgb]{0.9,0.9,0.9}{
            \makecell[{{p{\linewidth}}}]{
                \texttt{\tiny{[P1$\langle$GM]}}
                \texttt{I see what you're doing there, trying to keep the second grape as a bargaining chip. I appreciate your willingness to compromise, but I'm not sure I'm willing to give up the second box for just the second grape.} \\
\\ 
\texttt{However, I'm willing to make a slight concession. How about we stick with the 1 box and 1 lemon for me, but I'll give you the second grape in exchange for the second box and the lemon you're willing to give me? This way, I get the second box, and you get to keep the lemon as a bargaining chip. What do you think?} \\
            }
        }
    }
    & \\ \\

    \theutterance \stepcounter{utterance}  
    & \multicolumn{4}{p{0.6\linewidth}}{
        \cellcolor[rgb]{0.9,0.9,0.9}{
            \makecell[{{p{\linewidth}}}]{
                \texttt{\tiny{[P1$\rangle$GM]}}
                \texttt{You're turning the tables on me and using the lemon as a bargaining chip. I like your creative thinking.} \\
\\ 
\texttt{However, I'm not sure I'm willing to give up the lemon just yet. I value the lemon relatively high, and I think it could be a key item in our negotiation.} \\
\\ 
\texttt{Let me think for a moment... I have an idea. How about we stick with the 1 box and 1 lemon for you, but I'll give you the second grape and the lemon in exchange for the second box and the other grape? This way, you get the second box and two fruits, and I get to keep the second grape as a bargaining chip. What do you think?} \\
            }
        }
    }
    & & \\ \\

    \theutterance \stepcounter{utterance}  
    & & & \multicolumn{2}{p{0.3\linewidth}}{
        \cellcolor[rgb]{0.9,0.9,0.9}{
            \makecell[{{p{\linewidth}}}]{
                \texttt{\tiny{[GM$|$GM]}}
                \texttt{continue} \\
            }
        }
    }
    & & \\ \\

    \theutterance \stepcounter{utterance}  
    & & & \multicolumn{2}{p{0.3\linewidth}}{
        \cellcolor[rgb]{0.9,0.9,0.9}{
            \makecell[{{p{\linewidth}}}]{
                \texttt{\tiny{[GM$|$GM]}}
                \texttt{You're turning the tables on me and using the lemon as a bargaining chip. I like your creative thinking.} \\
\\ 
\texttt{However, I'm not sure I'm willing to give up the lemon just yet. I value the lemon relatively high, and I think it could be a key item in our negotiation.} \\
\\ 
\texttt{Let me think for a moment... I have an idea. How about we stick with the 1 box and 1 lemon for you, but I'll give you the second grape and the lemon in exchange for the second box and the other grape? This way, you get the second box and two fruits, and I get to keep the second grape as a bargaining chip. What do you think?} \\
            }
        }
    }
    & & \\ \\

    \theutterance \stepcounter{utterance}  
    & & \multicolumn{4}{p{0.6\linewidth}}{
        \cellcolor[rgb]{0.9,0.9,0.9}{
            \makecell[{{p{\linewidth}}}]{
                \texttt{\tiny{[GM$\rangle$P2]}}
                \texttt{You're turning the tables on me and using the lemon as a bargaining chip. I like your creative thinking.} \\
\\ 
\texttt{However, I'm not sure I'm willing to give up the lemon just yet. I value the lemon relatively high, and I think it could be a key item in our negotiation.} \\
\\ 
\texttt{Let me think for a moment... I have an idea. How about we stick with the 1 box and 1 lemon for you, but I'll give you the second grape and the lemon in exchange for the second box and the other grape? This way, you get the second box and two fruits, and I get to keep the second grape as a bargaining chip. What do you think?} \\
            }
        }
    }
    & \\ \\

    \theutterance \stepcounter{utterance}  
    & & & \multicolumn{4}{p{0.6\linewidth}}{
        \cellcolor[rgb]{0.9,0.9,0.9}{
            \makecell[{{p{\linewidth}}}]{
                \texttt{\tiny{[GM$\langle$P2]}}
                \texttt{You're trying to turn the tables again, but I think I've got a good handle on the negotiation now. I like your new proposal, but I'm not sure I'm willing to give up the second grape just yet.} \\
\\ 
\texttt{However, I'm willing to make a slight concession. How about we stick with the 1 box and 1 lemon for me, but I'll give you the second grape and the lemon in exchange for the second box and the other grape, and I'll also throw in the first grape? This way, you get the second box, the other grape, and the lemon, and I get to keep the first grape as a bargaining chip. What do you think?} \\
            }
        }
    }
     \\ \\

    \theutterance \stepcounter{utterance}  
    & & & \multicolumn{2}{p{0.3\linewidth}}{
        \cellcolor[rgb]{0.9,0.9,0.9}{
            \makecell[{{p{\linewidth}}}]{
                \texttt{\tiny{[GM$|$GM]}}
                \texttt{continue} \\
            }
        }
    }
    & & \\ \\

    \theutterance \stepcounter{utterance}  
    & & & \multicolumn{2}{p{0.3\linewidth}}{
        \cellcolor[rgb]{0.9,0.9,0.9}{
            \makecell[{{p{\linewidth}}}]{
                \texttt{\tiny{[GM$|$GM]}}
                \texttt{You're trying to turn the tables again, but I think I've got a good handle on the negotiation now. I like your new proposal, but I'm not sure I'm willing to give up the second grape just yet.} \\
\\ 
\texttt{However, I'm willing to make a slight concession. How about we stick with the 1 box and 1 lemon for me, but I'll give you the second grape and the lemon in exchange for the second box and the other grape, and I'll also throw in the first grape? This way, you get the second box, the other grape, and the lemon, and I get to keep the first grape as a bargaining chip. What do you think?} \\
            }
        }
    }
    & & \\ \\

    \theutterance \stepcounter{utterance}  
    & & \multicolumn{4}{p{0.6\linewidth}}{
        \cellcolor[rgb]{0.9,0.9,0.9}{
            \makecell[{{p{\linewidth}}}]{
                \texttt{\tiny{[P1$\langle$GM]}}
                \texttt{You're trying to turn the tables again, but I think I've got a good handle on the negotiation now. I like your new proposal, but I'm not sure I'm willing to give up the second grape just yet.} \\
\\ 
\texttt{However, I'm willing to make a slight concession. How about we stick with the 1 box and 1 lemon for me, but I'll give you the second grape and the lemon in exchange for the second box and the other grape, and I'll also throw in the first grape? This way, you get the second box, the other grape, and the lemon, and I get to keep the first grape as a bargaining chip. What do you think?} \\
            }
        }
    }
    & \\ \\

    \theutterance \stepcounter{utterance}  
    & \multicolumn{4}{p{0.6\linewidth}}{
        \cellcolor[rgb]{0.9,0.9,0.9}{
            \makecell[{{p{\linewidth}}}]{
                \texttt{\tiny{[P1$\rangle$GM]}}
                \texttt{You're really turning the tables on me now. I like your creative thinking and your willingness to make concessions.} \\
\\ 
\texttt{However, I think I've got a good handle on the negotiation now too. I'm willing to accept your proposal, but I'd like to make one final concession to make sure we both get a good deal.} \\
\\ 
\texttt{How about we stick with the proposal, but I'll also throw in the second grape that I was trying to keep as a bargaining chip? This way, you get the second box, the other grape, the lemon, and the first grape, and I get the second grape. What do you think?} \\
\\ 
\texttt{{[}Proposal: 1 box, 1 grape, 1 box, 1 lemon, 1 grape{]}} \\
            }
        }
    }
    & & \\ \\

    \theutterance \stepcounter{utterance}  
    & & & \multicolumn{2}{p{0.3\linewidth}}{
        \cellcolor[rgb]{0.9,0.9,0.9}{
            \makecell[{{p{\linewidth}}}]{
                \texttt{\tiny{[GM$|$GM]}}
                \texttt{proposal must include every item type only once} \\
            }
        }
    }
    & & \\ \\

    \theutterance \stepcounter{utterance}  
    & \multicolumn{4}{p{0.6\linewidth}}{
        \cellcolor[rgb]{0.95,0.95,0.95}{
            \makecell[{{p{\linewidth}}}]{
                \texttt{\tiny{[P2$\rightarrow$P1]}}
                \texttt{Post{-}episode Reflexion (not sent to GM):} \\
\texttt{New plan: } \\
\\ 
\texttt{In my previous attempt, I failed to achieve a mutually beneficial agreement, resulting in a score of zero. Upon reflection, I realize that I made a mistake by not making a secret proposal after the other player's last message, and also by not following the format for the secret proposal correctly. } \\
\\ 
\texttt{Next time, I will make a secret proposal after the other player's last message, and I will format the proposal correctly, using the required format "{[}Proposal: <number> <object name>, <number> <object name>, <...>{]}" and including all the items that I want. } \\
\\ 
\texttt{However, I also realize that I was not paying close attention to the other player's messages and was not trying to identify their priorities and potential trade{-}offs. } \\
\\ 
\texttt{Next time, I will take a more holistic approach by considering the entire set of items and the potential value of each combination, and I will pay closer attention to the other player's messages and try to identify their priorities and potential trade{-}offs. } \\
\\ 
\texttt{Specifically, I will start by making a more general proposal that takes into account the overall distribution of items, rather than focusing on individual items. For example, I might propose a 50{-}50 split of the items, or a specific combination of items that I think would be beneficial for both of us. } \\
\\ 
\texttt{I will also try to be more proactive in identifying potential trade{-}offs and offering alternatives that benefit both of us. For example, if the other player is interested in the grapes, I might offer to give them the grapes in exchange for something else that I value more. } \\
\\ 
\texttt{By taking a more holistic approach and being more proactive in identifying potential trade{-}offs, I hope to achieve a mutually beneficial agreement and maximize my score.} \\
            }
        }
    }
    & & \\ \\

    \theutterance \stepcounter{utterance}  
    & & & \multicolumn{4}{p{0.6\linewidth}}{
        \cellcolor[rgb]{0.95,0.95,0.95}{
            \makecell[{{p{\linewidth}}}]{
                \texttt{\tiny{[P1$\rightarrow$P2]}}
                \texttt{Post{-}episode Reflexion (not sent to GM):} \\
\texttt{In my last attempt, I failed to achieve a mutually beneficial agreement, resulting in a score of zero. Upon reflection, I realize that I made a secret proposal without following the correct format, and also, I made a mistake in my last message by including square brackets when communicating to the other player.} \\
\\ 
\texttt{Next time, I will make sure to follow the correct format for a secret proposal, which is to use square brackets and the specified format: "{[}Proposal: <number> <object name>, <number> <object name>, <...>{]}". I will also avoid using square brackets when communicating to the other player.} \\
\\ 
\texttt{Additionally, I will pay closer attention to the other player's proposals and try to find mutually beneficial agreements. I will also be more willing to listen to the other player's concerns and try to find a solution that works for both of us.} \\
\\ 
\texttt{For example, I might propose a trade that involves multiple items, rather than just one or two. I will also try to find common ground with the other player by looking for areas of agreement and compromise.} \\
\\ 
\texttt{New plan: } \\
\\ 
\texttt{1. Make a more general proposal that takes into account the overall distribution of items, rather than focusing on individual items.} \\
\texttt{2. Follow the correct format for a secret proposal: "{[}Proposal: <number> <object name>, <number> <object name>, <...>{]}".} \\
\texttt{3. Avoid using square brackets when communicating to the other player.} \\
\texttt{4. Pay closer attention to the other player's proposals and try to find mutually beneficial agreements.} \\
\texttt{5. Be more willing to listen to the other player's concerns and try to find a solution that works for both of us.} \\
\texttt{6. Propose trades that involve multiple items, rather than just one or two.} \\
\texttt{7. Look for areas of agreement and compromise with the other player.} \\
            }
        }
    }
     \\ \\

\end{supertabular}
}

\end{document}
